\documentclass[a4paper,10pt]{article}
\usepackage[left=0.9in,right=0.9in,top=0.8in,bottom=0.8in]{geometry}
\usepackage[utf8]{inputenc}
\usepackage{hyperref}
\usepackage{xcolor}

\definecolor{darkblue}{RGB}{0,119,181}
\definecolor{linkedinblue}{RGB}{0,119,181}
\hypersetup{colorlinks=true,linkcolor=linkedinblue,urlcolor=linkedinblue,pdfborder={0 0 0}}

\setlength{\parindent}{0pt}
\setlength{\parskip}{8pt}

\begin{document}
\pagestyle{empty}

{\color{darkblue}
\textbf{Coretura AB}\\
AI Lab\\
Gothenburg, Sweden
}

\vspace{14pt}

\textbf{Dear Hiring Team,}


I am applying for the AI Platform Engineer role in Coretura's AI Lab. I work in the Software Factory at Volvo Group Digital Technology \& Operations (DTO), where I build and run an on-premises AI platform and the agents on top of it, so I already know commercial-vehicle software, its security constraints and how one of your two owners works. \textbf{I am available from 1 January 2027.}


What draws me to Coretura is the AI Lab's mandate: one shared platform that every team uses to build and run AI products safely and reliably. That is what I have been doing at Volvo Group, on a smaller scale. This year I built the first RAG system in DTO, a support chatbot that two teams use every day and that cut ticket processing time by 35\%. I followed it with DevBot, which cut senior pull-request review time from 6 hours to 60 minutes, and OfficeBot, which saves about 90\% of the time it takes to build a slide deck. All of them run on a local model platform that I set up and operate, with no data going to an outside AI service. At Volvo Group HackDay 2026 I led a newly formed team to \textbf{2nd place of 16 teams} with ClaimGenerationBot.


Your job description asks for someone who sets the platform's direction and also takes on the hands-on work: model access across providers, retrieval, agent harnesses, security, cost, and the observability that shows whether an AI product gives the right answer and not just whether it is running. That matches how I work. My models sit behind OpenAI-compatible endpoints with failover, so a product changes model without code changes. Answers carry sources and a confidence score, and low-confidence answers go to a specialist. Knowledge bases are versioned separately from the application, so a change is reviewed and traceable. I set up Claude Code for my team through Volvo's AI gateway with short-lived tokens instead of API keys, and I use AI coding agents in my own work every day.


The project closest to Coretura's mission is \textbf{Volvo Smart Vehicle}: three AI agents working together on one truck record. Co-Driver talks with the driver in the cab, ClaimBot settles who pays for a repair, and Fleet Manager runs the haulier's office. They hand work to each other in real time, so a warning lamp on the dashboard becomes a booked workshop slot and a settled claim before the truck arrives, in about a second per voice turn. I believe an intelligent cabin like this belongs on a shared software platform for commercial vehicles, and I would bring the experience of having built it end to end.


You are also looking for someone who builds to discover and then shares what they learn. I have presented twice to the whole of DTO, with more than 70 colleagues registered each time, I run knowledge-sharing sessions in my team, and I wrote the team's coding handbook, because the biggest multiplier for any engineer is the team around them.


I would welcome the chance to discuss how I can contribute to the AI Lab, and to show a 10-minute demo of Smart Vehicle.

\vspace{6pt}

Sincerely,\\[10pt]
\textbf{Harvad (Hongzhi) Li}

\vspace{10pt}
{\color{darkblue}\rule{\linewidth}{0.6pt}}

{\color{darkblue}\small
Gothenburg, Sweden \textbar{} \href{mailto:hongzhili01@gmail.com}{hongzhili01@gmail.com} \textbar{} +46 76 006 7977\\
LinkedIn: \href{https://www.linkedin.com/in/harvad/}{linkedin.com/in/harvad} \textbar{} GitHub: \href{https://github.com/bluehawana}{github.com/bluehawana}
}

\end{document}
